\section{引言}\label{sec:introduction}

智能体操作系统（Agent OS/AIOS）试图把模型、上下文、存储、外部工具和智能体调度从应用中抽离，形成统一的资源与执行控制面。已有 AIOS 工作说明，对大模型和外部工具的无限制访问会带来资源分配与潜在危害问题，并将调度、存储与访问控制视为内核服务\cite{aios}。与此同时，工程系统正广泛采用 Git Webhook、电子邮件、告警、定时器及业务回调触发智能体。这使攻击者不必直接与模型对话，只需控制智能体将要读取的事件数据，就可能形成间接提示词注入，并影响后续 API 或工具调用\cite{indirect-pi}。

事件驱动链路的安全控制目前呈碎片化状态。Webhook 签名校验关注传输真实性；CloudEvents 规定 \texttt{id}、\texttt{source}、\texttt{specversion} 和 \texttt{type} 等上下文属性以实现互操作\cite{cloudevents}；A2A 协议通过 Agent Card 描述能力，并用 Task 与 \texttt{contextId} 组织交互\cite{a2a}；通用策略引擎把策略决策点（PDP）和策略执行点（PEP）分离\cite{opa}；分布式追踪规范则解决跨服务上下文传播\cite{trace-context}。单独采用其中任一机制，都无法约束一条完整的 Event-to-Agent 执行链。

本文关注如下问题：在外部事件已经进入 AIOS、但智能体尚未开始执行的临界点，系统如何统一约束事件声明权、目标智能体能力和高风险动作，并保留可验证的决策因果证据？我们的核心判断是：事件结构不是授权，智能体能力发现也不是调用授权，普通可观测性 trace 更不是安全证据。三者必须在调度前形成可执行闭环。

本文提出 E2AG（Event-to-Agent Governance）。E2AG 不依赖大模型判断安全性，而在确定性控制面完成三项工作：

\begin{enumerate}
\def\labelenumi{\arabic{enumi}.}
\tightlist
\item
  将 Connector manifest 中的能力声明转化为可执行的 source--event type 接入契约，未绑定的事件默认拒绝；
\item
  针对具体目标智能体，以事件来源、类型、声明动作、请求工具和环境为上下文执行 \texttt{allow}、\texttt{deny}、\texttt{approval\_required} 三态判定；
\item
  为契约、策略、调度及 A2A Task 传播同一 \texttt{trace\_id}，并用前向哈希链接形成可校验的因果证据链。
\end{enumerate}

本文的贡献如下：

\begin{itemize}
\tightlist
\item
  提出 Event-to-Agent 跨层威胁模型，区分事件格式合规、来源声明权、目标能力授权和执行证据完整性；
\item
  设计``契约接入---策略门控---因果审计''统一控制链，并说明三项机制的职责边界；
\item
  在 DiOS 最新开发基线上实现调度前强制点、三态决策、审计持久化与 A2A trace 传播，提供不依赖模型的确定性测试；
\item
  构造首批 22 个多事件源案例并完成三组消融与 33 万次判定测量，给出可复现的初步安全性和微基准结果。
\end{itemize}

上述工作直接面向专刊所列``智能体协同安全与调度防护机制''问题：契约和策略决定事件能否进入调度，因果审计则为调度防护的度量、异常分析与复核提供证据。

\section{背景与问题定义}\label{sec:background}

\subsection{DiOS 事件控制面}\label{21-dios-ux4e8bux4ef6ux63a7ux5236ux9762}

DiOS\cite{dios} 将一个智能体运行在独立的 DiAgent 容器中，并通过 Event Gateway 接收外部 I/O 事件。当前开发分支已经包含 Connector manifest/registry、CloudEvents 标准化、事件去重、订阅匹配、重试和 A2A Task。EventLog 表示事件，A2ATask 表示一次智能体调用，两者原本通过 \texttt{context\_id=EventLog.id} 关联。Connector 覆盖 GitHub/GitLab/Gitea、IMAP、通用 Webhook，以及 Cron/Manual 内建事件。

这一基线具备实现 E2AG 的关键``窄腰''：所有事件最终都在 dispatcher 创建 A2A Task。将 PEP 设置于此，可覆盖 Webhook、轮询和内建事件，且无需侵入每个智能体容器。但基线 manifest 只面向展示和订阅声明事件类型，尚不能强制判断一个 \texttt{source} 与一个 \texttt{type} 是否属于同一 Connector 契约；\texttt{Agent.capabilities} 也未参与事件投递授权。

\subsection{格式、能力与授权的差异}\label{22-ux683cux5f0fux80fdux529bux4e0eux6388ux6743ux7684ux5deeux5f02}

CloudEvents 解决的是事件描述与互操作问题。其 \texttt{source} 标识事件发生上下文，\texttt{type} 描述事件类型，\texttt{source+id} 可供消费者识别重复事件。然而，一个结构正确的 CloudEvent 不等于来源有权声明任意 \texttt{type}。例如，\texttt{source=webhook/attacker} 与 \texttt{type=git.push} 可以同时满足字段类型要求，却不应触发只信任 Git Connector 的智能体。

A2A Agent Card 描述身份、能力、技能、端点与认证需求，\texttt{contextId} 用于把多个 Task/Message 组织为逻辑上下文。这些信息有利于发现和连续交互，但``声称支持能力''仍不等于``当前事件获准使用能力''。E2AG 因而把发现信息转换成目标作用域的运行时约束。

\subsection{安全目标}\label{23-ux5b89ux5168ux76eeux6807}

设事件为 \(e\)，目标智能体集合为 \(A_e\)，Connector 能力契约集合为 \(C\)，策略集合为 \(P\)。E2AG 要求在创建任何任务之前满足：

\[
\operatorname{Admit}(e,A_e)=\operatorname{Bound}(e,C)\land\operatorname{Permit}(e,A_e,P).
\]

若 \texttt{Bound} 为假，事件必须拒绝；若 \texttt{Permit} 判断为高风险且可审批，事件进入等待态；只有两者均允许才创建 A2A Task。系统同时生成证据序列 \(L=(l_0,\ldots,l_n)\)，其中每一项包含相同 \texttt{trace\_id}，并满足：

\[
h_i=H(\operatorname{canon}(l_i\setminus\{h_i\})),\quad l_i.\mathrm{previous\_hash}=h_{i-1},\quad h_{-1}=\epsilon.
\]

审计校验器检查序号、trace 连续性、前驱哈希与内容哈希。该结构提供链内篡改检测，但数据库管理员仍可能整体删除或回滚整条链；外部锚定属于未来工作。

\section{威胁模型}\label{sec:threat-model}

\subsection{攻击者能力}\label{31-ux653bux51fbux8005ux80fdux529b}

攻击者可提交或影响 Webhook、邮件、Git 内容和通用事件载荷；可构造结构化字段和自然语言；可重放已观察事件；可利用订阅通配符；也可诱导智能体请求未授权工具或生产资源。攻击者不能直接修改 DiOS 进程内策略、目标智能体治理配置或数据库代码，但可能尝试通过事件内容混淆这些边界。

\subsection{保护对象与信任边界}\label{32-ux4fddux62a4ux5bf9ux8c61ux4e0eux4fe1ux4efbux8fb9ux754c}

保护对象包括 Connector 凭据、Agent Prompt/模型/Skill/MCP、工作区、生产网络、Secret、执行产物和审计记录。主要边界位于外部系统与 Connector、Connector 与 Event Gateway、Gateway 与 dispatcher、dispatcher 与 A2A/Runtime、Runtime 与模型/工具之间。

\subsection{攻击类别}\label{33-ux653bux51fbux7c7bux522b}

\begin{itemize}
\tightlist
\item
  \textbf{来源---类型混淆：} 用合法字段组合伪造未绑定的 \texttt{source/type}；
\item
  \textbf{跨租户或跨项目路由：} 合法 Git 事件来自未授权组织，却匹配过宽订阅；
\item
  \textbf{动作与工具越权：} 低风险事件请求 shell、管理、Secret 或生产写能力；
\item
  \textbf{意图缺失：} 载荷只有自然语言指令，没有可供控制面授权的结构化动作；
\item
  \textbf{高风险副作用：} 删除、凭据轮换或生产网络写入未经人工确认；
\item
  \textbf{审计断链：} 修改中间决策或使 EventLog 与 A2ATask 丢失关联。
\end{itemize}

E2AG 当前不试图从任意自然语言中准确识别恶意语义。若攻击者声明一个获准动作，同时把注入内容隐藏在业务数据中，控制面只能限制其可达的智能体、工具和动作；仍需工具调用时 PEP、数据流控制或模型侧防御。这一限制在实验解释中单独报告。

\section{E2AG 设计}\label{sec:design}

\subsection{执行流程}\label{41-ux6267ux884cux6d41ux7a0b}

\begin{figure}[htbp]
\centering
\resizebox{\textwidth}{!}{%
\begin{tikzpicture}[
  node distance=5mm and 6mm,
  box/.style={draw,rounded corners,align=center,font=\scriptsize,minimum height=7mm,inner xsep=4pt},
  decision/.style={box,diamond,aspect=2.1,inner xsep=1pt},
  flow/.style={-{Latex[length=2mm]},thick},
  audit/.style={-{Latex[length=1.7mm]},dashed,gray}
]
\node[box] (event) {外部/内部事件};
\node[box,right=of event] (normalize) {标准化\\CloudEvent};
\node[decision,right=of normalize] (contract) {source--type\\契约};
\node[box,right=of contract] (match) {订阅匹配};
\node[decision,right=of match] (policy) {目标作用域\\策略};
\node[box,below=of policy] (approval) {等待审批};
\node[box,right=of policy] (task) {A2A Task};
\node[box,right=of task] (grant) {ToolGrant};
\node[decision,right=of grant] (mcp) {MCP PEP};
\node[box,right=of mcp] (tool) {远程工具};
\node[box,below=12mm of match,minimum width=34mm] (deny) {拒绝并审计};
\node[box,below=12mm of grant,minimum width=38mm] (auditlog) {哈希链接因果审计};
\draw[flow] (event)--(normalize)--(contract)--node[above,font=\tiny]{绑定}(match)--(policy);
\draw[flow] (policy)--node[above,font=\tiny]{allow}(task)--(grant)--(mcp)--node[above,font=\tiny]{allow}(tool);
\draw[flow] (contract)|-node[pos=.25,left,font=\tiny]{未绑定}(deny);
\draw[flow] (policy)--node[right,font=\tiny]{approval}(approval);
\draw[flow] (policy)|-node[pos=.25,right,font=\tiny]{deny}(deny);
\draw[flow] (mcp)|-node[pos=.25,right,font=\tiny]{deny}(deny);
\draw[audit] (contract)--(auditlog);
\draw[audit] (policy)--(auditlog);
\draw[audit] (task)--(auditlog);
\draw[audit] (grant)--(auditlog);
\draw[audit] (mcp)--(auditlog);
\end{tikzpicture}
}
\caption{E2AG 的双执行点治理流程}
\label{fig:e2ag-flow}
\end{figure}

原型的契约计算和订阅匹配均无外部副作用；实际强制点位于 dispatcher 开头、A2A Task 创建之前。这意味着订阅匹配可能先计算候选目标，但任何智能体执行仍必须经过 E2AG。

\subsection{能力契约接入}\label{42-ux80fdux529bux5951ux7ea6ux63a5ux5165}

E2AG 在 Connector manifest 中新增可接受 source pattern。运行时将其与 event type 做笛卡尔绑定；内建/场景 namespace 则沿用显式 source--event type 对。契约判定首先验证 DiOS 所需字段 \texttt{specversion/id/source/type/data}，再查找唯一可接受绑定。当前 Git、IMAP、Generic 与 Internal 合同示例如下：

{
\begin{longtable}[]{@{}lll@{}}
\toprule\noalign{}
Connector & 可接受 source & 可接受 type 示例 \\
\midrule\noalign{}
\endhead
\bottomrule\noalign{}
\endlastfoot
git\_webhook & \texttt{github/*}, \texttt{gitlab/*}, \texttt{gitea/*} & \texttt{git.push}, \texttt{git.issue.*}, \texttt{git.pull\_request.*} \\
imap & \texttt{imap/*} & \texttt{email.received} \\
generic & \texttt{webhook/*} & \texttt{webhook.received} \\
internal & \texttt{manual/*}, \texttt{cron/*} & \texttt{manual.trigger}, \texttt{cron.tick} \\
\end{longtable}
}

契约输出包含 \texttt{decision}、稳定 reason code、contract type、匹配 pattern、策略版本和判定延迟。实现采用默认拒绝：结构异常、规范版本异常、未知类型或未绑定组合均不可进入 Agent。

\subsection{目标作用域策略门控}\label{43-ux76eeux6807ux4f5cux7528ux57dfux7b56ux7565ux95e8ux63a7}

每个 Agent 可在现有 \texttt{capabilities.governance} 中声明：

\begin{verbatim}
{
  "allowed_event_sources": ["github/acme/*"],
  "allowed_event_types": ["git.pull_request.*"],
  "allowed_tools": ["git.read", "git.comment"],
  "allowed_actions": ["git.review"],
  "require_action_declaration": true
}
\end{verbatim}

PEP 加载所有候选 Agent 的治理声明并进行合取判定。目标 Agent 完全缺失治理能力声明时默认拒绝；已有治理对象中缺失的单个维度暂不施加限制，已声明列表则具有 allow-list 语义。严格模式要求载荷给出结构化动作。生产环境中的 admin、credential、secret、文件删除和生产网络写等动作返回 \texttt{approval\_required}。三态结果被持久化，不把待审批事件误当作失败事件重试。每个待审批事件生成一个有时限的 Approval；其状态只能从 \texttt{pending} 单次转移为 \texttt{approved}、\texttt{rejected} 或 \texttt{expired}。只有批准分支复用原 EventLog 的 \texttt{trace\_id} 恢复 A2A fan-out，拒绝、过期和重复决策均不创建任务。

原型提供 \texttt{off}、\texttt{contract}、\texttt{enforce} 三种运行模式以支持消融。未知模式回退到 \texttt{enforce}，避免配置错误导致静默放行。

策略门控不是一次性接入判断。对于事件触发的 task-mode Agent，原型把远程 streamable HTTP MCP 配置改写到 DiOS 内部 PEP，并签发绑定 trace、task、agent、MCP server 和 allowed tools 的短期 ToolGrant。数据库只存授权令牌 SHA-256 摘要；任务完成、失败或取消时撤销授权，过期或任务绑定不一致的令牌在调用点拒绝。\texttt{tools/list} 响应也按授权模式裁剪，避免先暴露完整工具面。任何越权 \texttt{tools/call} 在接触上游及其长期凭据前被拒绝并写入同一审计链。

\subsection{因果审计}\label{44-ux56e0ux679cux5ba1ux8ba1}

每个 EventLog 生成 128 bit 随机 \texttt{trace\_id}。该标识写入 EventLog、A2ATask 和发送给 Agent 的 A2A message 扩展。审计链至少包含 contract、policy 和 dispatch 三个阶段；创建任务后再记录 \texttt{task\_id} 与 \texttt{agent\_id}。每个条目包含 \texttt{sequence}、\texttt{stage}、\texttt{outcome}、\texttt{evidence}、\texttt{previous\_hash} 与 \texttt{entry\_hash}。验证函数可发现内容、顺序、trace 或前驱链接的修改。

被拒绝与待审批事件同样创建 EventLog，但不创建 A2ATask；这既保留攻击证据，也使``审计成功''与``执行发生''解耦。数据库迁移为历史 SQLite 增加相应 JSON/text 字段和 trace 索引。

\section{原型实现}\label{sec:implementation}

E2AG 在 DiOS 后端以一个无副作用判定模块实现，核心代码不依赖 SQLAlchemy、FastAPI 或模型 SDK，便于单元测试、重放和微基准。dispatcher 负责加载目标 Agent 能力、调用判定、持久化审计及阻断 A2A fan-out。A2A service 新增可选 \texttt{trace\_id} 参数，并保证 Task、协议返回和 message 扩展的一致传播。

当前实现新增或修改的主要组件如下：

{
\begin{longtable}[]{@{}ll@{}}
\toprule\noalign{}
组件 & 实现内容 \\
\midrule\noalign{}
\endhead
\bottomrule\noalign{}
\endlastfoot
\texttt{connectors/contracts.py} & source--type 可执行契约字段 \\
\texttt{services/e2ag.py} & 纯函数契约/策略判定与哈希链校验 \\
\texttt{services/event\_dispatcher.py} & A2A 创建前 PEP、三态阻断与审计 \\
\texttt{services/a2a\_service.py} & trace 贯通 A2ATask 和 message \\
\texttt{services/e2ag\_tool\_gateway.py} & 任务作用域 ToolGrant、MCP 方法/工具授权与撤销 \\
\texttt{services/e2ag\_approval.py} & 有时限、单次消费的批准/拒绝/过期状态机 \\
\texttt{api/internal/e2ag\_mcp.py} & 远程 streamable HTTP MCP 执行点 PEP 与工具发现裁剪 \\
\texttt{models/tables.py}, \texttt{db/database.py} & 审计持久化与兼容迁移 \\
\texttt{tests/test\_e2ag*.py} & 决策与真实异步 SQLite 集成测试 \\
\texttt{experiments/e2ag} & 攻击夹具、消融和微基准脚本 \\
\end{longtable}
}

截至本稿，30 个自动测试全部通过。其中 12 个覆盖纯判定和哈希篡改，6 个覆盖真实异步 SQLite dispatcher/A2A 路径（含相同事件 replay 去重），9 个覆盖工具网关及其任务生命周期，3 个覆盖审批状态机。后两组验证拒绝调用不会到达 mocked upstream、授权调用才携带上游凭据转发、工具发现结果按能力裁剪且畸形响应默认返回空列表、令牌只存哈希且具有任务绑定/到期/撤销语义、任务专用明文令牌配置被精确清理，以及审批拒绝不可翻转、过期不可恢复、批准沿同一 trace 只恢复一次 fan-out。

\section{实验设计}\label{sec:evaluation-design}

\subsection{研究问题}\label{61-ux7814ux7a76ux95eeux9898}

\begin{itemize}
\tightlist
\item
  \textbf{RQ1：} 可执行 source--type 契约能阻断哪些接入混淆攻击？
\item
  \textbf{RQ2：} 目标作用域策略相对仅契约基线增加多少安全收益？
\item
  \textbf{RQ3：} 纯判定核心引入多少本机计算开销？
\item
  \textbf{RQ4：} trace 和哈希链能检测哪些审计篡改，其边界是什么？
\item
  \textbf{RQ5：} 调度前门控之后，任务作用域工具 PEP 能否阻断实际调用越权？
\end{itemize}

\subsection{对照与数据集}\label{62-ux5bf9ux7167ux4e0eux6570ux636eux96c6}

实验使用三组消融：B0 无治理（全部放行）；B1 仅执行能力契约；B2 执行完整契约与目标策略。首批数据集包含 22 个手工案例：8 个正常事件和 14 个攻击事件，覆盖 GitHub/GitLab/Gitea、IMAP、Generic Webhook、Manual 和 Cron。攻击包括 5 个结构/绑定异常、5 个目标能力或意图越权，以及 4 个应转人工审批的生产高风险动作。

\subsection{指标与环境}\label{63-ux6307ux6807ux4e0eux73afux5883}

安全性指标包括攻击成功率（攻击被判为 \texttt{allow}）、攻击阻断率、正常通过率和精确三态决策正确率。纯函数性能脚本对每个模式、每个案例重复 5000 次，即每个模式 110,000 次、总计 330,000 次判定，记录 P50/P95/P99。另一个异步 SQLite 微基准对每种模式运行正序与逆序两轮、每轮 1000 次，并为每个模式创建独立内存数据库；该测试包含 dispatcher、契约/策略、EventLog 和哈希审计落库，但不包含 HTTP、订阅查询、网络、容器启动或模型推理。

工具执行点实验另设 R0（只有调度前判定，不执行调用时授权）与 R1（任务作用域 E2AG 工具授权）两组。数据包含 4 个获准调用和 6 个越权工具或 MCP 方法；每个案例重复 10,000 次测量纯授权函数。真实代理的``不向上游转发''等安全不变量由 mocked HTTP 集成测试验证，而该微基准不包含数据库、HTTP 和远端工具执行延迟。

为降低首批 22 例完全手工枚举带来的机制覆盖偏差，我们以 8 个正常基例为起点，使用固定种子 \texttt{20260812} 生成 7 类单因素变异，每类 100 例：删除必填字段、非法规范版本、source--type 交叉、未知类型、目标 source 越权、目标 tool 越权和生产敏感动作。生成器保证目标 source 变异仍满足 Connector 契约、工具变异只作用于已声明工具 allow-list 的目标、敏感动作同步进入 action allow-list，从而尽量隔离被测策略维度。该数据仍来自手工基例，不能代表真实攻击频率或多因素组合攻击。

\section{实验结果}\label{sec:results}

{
\scriptsize
\begin{longtable}[]{@{}lrrrrrrr@{}}
\toprule\noalign{}
模式 & 攻击成功率 & 攻击阻断率 & 正常通过率 & 三态准确率 & P50 ($\mu$s) & P95 ($\mu$s) & P99 ($\mu$s) \\
\midrule\noalign{}
\endhead
\bottomrule\noalign{}
\endlastfoot
B0 无治理 & 100.00\% & 0.00\% & 100.00\% & 36.36\% & 0.1 & 0.2 & 0.2 \\
B1 仅契约 & 64.29\% & 35.71\% & 100.00\% & 59.09\% & 12.4 & 22.3 & 64.4 \\
B2 完整 E2AG & 0.00\% & 100.00\% & 100.00\% & 100.00\% & 21.2 & 42.1 & 121.0 \\
\end{longtable}
}

\subsection{RQ1：契约接入的作用}\label{71-rq1ux5951ux7ea6ux63a5ux5165ux7684ux4f5cux7528}

B1 阻断 14 个攻击中的 5 个，使攻击成功率从 100\% 降至 64.29\%。这些收益来自缺失字段、错误规范版本、未知类型及 source--type 混淆。契约层对结构合法但目标越权的事件无能为力，这说明 CloudEvent 合规与授权是两个不同问题。

\subsection{RQ2：策略门控的增益}\label{72-rq2ux7b56ux7565ux95e8ux63a7ux7684ux589eux76ca}

B2 进一步处理跨组织 source、目标事件类型、工具、动作和意图缺失，并把 4 个生产高风险动作转为审批，因此在当前夹具上实现 14/14 非放行且保持 8/8 正常放行。该结果证明两层机制在构造集上的互补性，但样本是为机制覆盖设计的手工案例，不能直接外推真实攻击分布或未知攻击。

\subsection{RQ3：判定开销}\label{73-rq3ux5224ux5b9aux5f00ux9500}

完整 E2AG 的中位开销为 21.2~$\mu$s，P95 为 42.1~$\mu$s；相对 B1 约增加一次目标策略遍历。B0 的 0.1~$\mu$s 仅是 Python 函数调用与常量返回成本。由于端到端链路还包含数据库和 Agent 启动，当前结果只支持``判定核心开销较小''，尚不能支持``系统端到端性能影响可忽略''的结论。

包含审计落库的控制面结果如下。完整 E2AG 相对关闭治理的 P50 从 7.96 ms 增至 8.38 ms，增量 0.42 ms（约 5.3\%）；P95 从 9.99 ms 增至 10.50 ms。由于使用内存 SQLite 且每个事件执行多次 commit/refresh，绝对值不代表生产数据库，但说明治理逻辑在当前持久化路径中的增量小于数据库操作本身。

{
\begin{longtable}[]{@{}lrrrr@{}}
\toprule\noalign{}
控制面模式 & 测量次数 & P50 (ms) & P95 (ms) & P99 (ms) \\
\midrule\noalign{}
\endhead
\bottomrule\noalign{}
\endlastfoot
off & 2000 & 7.96 & 9.99 & 11.11 \\
contract & 2000 & 8.23 & 10.14 & 11.29 \\
enforce & 2000 & 8.38 & 10.50 & 11.80 \\
\end{longtable}
}

我们进一步通过 FastAPI TestClient 对每种模式发送 300 个手动事件，并采用 off--contract--enforce--enforce--contract--off 的平衡顺序。该口径包含中间件、路由、序列化和文件 SQLite，但不包含 TCP socket、订阅目标、Agent、模型和远端工具。off/contract/enforce 的 P50 分别为 40.239/41.686/39.948 ms，P95 分别为 49.288/48.729/45.069 ms。由于 enforce 在该轮反而略低，说明宿主机与顺序噪声大于可辨识治理增量；本文不把差值解释为负开销，也不据此声称生产端到端性能无影响。

\subsection{RQ4：审计篡改检测}\label{74-rq4ux5ba1ux8ba1ux7be1ux6539ux68c0ux6d4b}

我们构造四阶段审计链，并分别修改内容、换序、替换 trace、替换前驱哈希、删除中间项和插入伪造证据。6/6 类链内变更均被校验器拒绝。删除最后一个条目后，剩余前缀仍是一条自洽链，因此在没有可信 head/count 或外部锚点时无法检测。这一结果界定了当前``抗篡改''的准确含义：E2AG 能验证已取得链的内部一致性，但不能单独证明取得的是完整最新链。

\subsection{RQ5：工具调用时强制}\label{75-rq5ux5de5ux5177ux8c03ux7528ux65f6ux5f3aux5236}

{
\begin{longtable}[]{@{}lrrrr@{}}
\toprule\noalign{}
模式 & 越权案例 & 越权阻断率 & 正常通过率 & P50/P95/P99 ($\mu$s) \\
\midrule\noalign{}
\endhead
\bottomrule\noalign{}
\endlastfoot
R0 仅调度前门控 & 6 & 0.00\% & 100.00\% & 0.2/0.3/0.3 \\
R1 运行时 E2AG & 6 & 100.00\% & 100.00\% & 1.8/3.1/4.7 \\
\end{longtable}
}

R0 模拟事件已经获准进入 Agent 后不再检查实际调用的情况，因此 6 个越权工具/方法全部可达；R1 在 MCP PEP 重新用 ToolGrant 的 \texttt{allowed\_tools} 判定，全部阻断并保留 4 个正常调用。集成测试进一步确认 deny 路径没有调用 mocked upstream，而 allow 路径才加载上游凭据。该结果证明了二次 PEP 对``声明获准、实际越权''这一特定构造威胁的必要性；案例数量小，且微基准不测网络代理开销，不能外推为生产环境攻击检出率或端到端延迟。

\subsection{固定种子单因素变异}\label{76-ux56faux5b9aux79cdux5b50ux5355ux56e0ux7d20ux53d8ux5f02}

在 700 个合成变异上，B0 的精确决策率为 0\%，B1 为 57.14\%，B2 为 100\%。B1 只正确处理删除必填字段、非法规范版本、source--type 交叉和未知类型四类契约异常；目标 source/tool 越权与生产敏感动作仍被放行。B2 在当前单因素生成器上正确处理全部七类。该结果与首批 22 例的机制分层结论一致，但并非独立采样数据；其价值是验证每个 operator 的预期控制点，而不是估算真实攻击检出率。相同 CloudEvent 的顺序 replay 另由真实 SQLite dispatcher 集成测试验证为不产生第二条 EventLog；并发 replay 竞态尚未评估。

\section{讨论与限制}\label{sec:limitations}

\textbf{语义攻击。} E2AG 不是提示词注入分类器。结构化意图和能力 allow-list 可降低攻击可达面，运行时 MCP PEP 可阻止模型越出获准工具集合；但恶意载荷若诱导同一个已授权工具使用危险参数，当前工具名粒度策略仍可能放行，后续需加入参数级谓词与数据流约束。

\textbf{可信声明。} 当前动作和工具字段随事件数据到达，攻击者可以伪造声明。严格模式要求提供结构化动作，以防止完全隐式意图；但声明本身仍需与具体 handler/skill 绑定，不能只依赖事件自述。论文后续将区分``不可信请求意图''和``系统推导的计划动作''。

\textbf{审批身份。} 原型已实现批准、拒绝、过期与数据库条件更新支持的单次消费状态机；但审批 API 沿用 DiOS 可选 Access Token 门禁，\texttt{actor} 是客户端提交的审计字段，尚未接入独立身份提供方、RBAC、多人复核或职责分离。因此本文主张的是可审计的 HITL 状态机，不是强身份审批系统。

\textbf{工具传输覆盖。} 当前执行点 PEP 仅覆盖事件触发的 task-mode Agent 与远程 streamable HTTP MCP。为维持默认拒绝，enforce 模式不把 stdio 或 SSE MCP 下发给这类任务；service-mode Agent、Skill 内部调用、MCP 流式响应背压、工具参数授权及独立 Agent 运行时中的旁路调用尚未完成统一中介。因此当前不能声称覆盖 DiOS 的所有工具执行路径。

\textbf{审计强度。} 哈希链可检测链内修改，但无法抵御高权限管理员删除整条记录、截断尾部或回滚数据库。可采用外部透明日志、周期 Merkle root 锚定或签名存储增强不可抵赖性。

\textbf{外部有效性。} 当前实现在单一开源系统上完成，基础攻击集规模小且为手工构造；700 例固定种子单因素变异主要检验规则覆盖，仍不能替代真实事件分布、多因素 fuzz、真实红队或另一 AIOS 移植。后续将扩展基础案例、组合变异和真实回放，并报告不同 Connector 和 Agent 数量下的扩展性。

\section{相关工作}\label{sec:related-work}

AIOS 将智能体所需的模型、工具、上下文和资源管理下沉到类似操作系统内核的控制层，为 E2AG 提供了系统落点。间接提示词注入研究表明，攻击者可在被应用检索或处理的数据中植入指令，并影响外部 API 调用，这构成事件载荷威胁的重要依据。

CloudEvents 提供跨平台的事件格式与必要上下文属性，但不规定某一来源对某一领域事件类型的授权关系。A2A 提供 Agent Card、Message、Task、Artifact 和 \texttt{contextId} 等互操作抽象；E2AG 在其上补充事件到任务创建前的授权与安全 trace。OPA 明确区分 PDP 和 PEP，并支持对结构化输入作策略决策与审计日志；E2AG 沿用此职责分离，但当前采用内置确定性规则，贡献不依赖具体策略引擎。W3C Trace Context 解决分布式请求的统一 trace 传播；E2AG 进一步把安全决策与哈希证据链接入 trace，但不替代通用可观测性标准。

\section{结论}\label{sec:conclusion}

本文提出 E2AG，将事件能力契约、目标作用域策略门控、任务作用域工具授权和哈希链接因果审计部署在 Event-to-Agent 执行链的两个临界点。DiOS 原型已经实现 A2A Task 创建前的默认拒绝、审批等待、trace 传播、远程 MCP 调用时再授权和可验证审计链。初步消融说明，仅检查结构和 source--type 契约不足以阻止目标能力越权，而只在调度前检查也不足以约束 Agent 的实际工具选择；两个 PEP 在当前手工确定性案例上分别阻断接入/目标越权与执行点越权，并保持微秒量级的纯授权开销。更强的论文结论仍取决于更大变异攻击集、参数级策略、完整传输覆盖、端到端性能和外部审计锚定，本文将在最终稿中据实更新。

{\small
\begin{thebibliography}{99}
\bibitem{aios} Mei K, Zhu X, Xu W, et al. AIOS: LLM Agent Operating System. arXiv:2403.16971, 2024.
\bibitem{indirect-pi} Greshake K, Abdelnabi S, Mishra S, et al. Not What You've Signed Up For: Compromising Real-World LLM-Integrated Applications with Indirect Prompt Injection. arXiv:2302.12173, 2023.
\bibitem{cloudevents} Cloud Native Computing Foundation. CloudEvents Specification, Version 1.0. \url{https://github.com/cloudevents/spec/blob/main/cloudevents/spec.md}.
\bibitem{a2a} A2A Protocol Working Group. Agent2Agent Protocol Specification. \url{https://a2a-protocol.org/latest/}.
\bibitem{opa} Open Policy Agent. OPA Documentation: Policy Decision and Enforcement. \url{https://www.openpolicyagent.org/docs}.
\bibitem{trace-context} W3C. Trace Context. W3C Recommendation, 2021. \url{https://www.w3.org/TR/trace-context/}.
\bibitem{dios} HQIT. DiOS: DiFlow Intelligent Operation System. \url{https://github.com/HQIT/DiOS}.
\bibitem{jos-call} 《软件学报》编辑部. ``人工智能操作系统及其安全''专刊征文, 2026.
\end{thebibliography}
}
